\section{The bank and the transaction also determine eligibility}
\label{sec:bank-compliance}

An ordinary bond can have straightforward payment terms and still be unavailable
for a particular transaction. The client may be prohibited from acquiring it;
the bank may be prohibited from dealing in it; or the bank may lack information
needed to establish which restrictions apply. The CoCo example therefore exposes
a general question. After explaining what an instrument promises, the system
must establish what the participating institutions and clients may do with it.
Changing the instrument from an AT1 bond to a preferred share, a fund or an
ordinary equity does not remove that question.

\subsection{Identify the institution before selecting its rules}

Consider a private-banking service provided through a Hong Kong office.
JPMorgan's public Asia disclosure identifies JPMorgan Chase Bank, N.A., Hong
Kong or Singapore branch as the provider of dealing, advisory, discretionary
management, banking and custody services, as notified to the client. It also
describes the Hong Kong branch's local regulatory supervision. The same page
identifies other legal entities in other jurisdictions. Consequently, the
J.P. Morgan brand alone cannot identify the legal person undertaking a trade;
the account and service documents must establish that person and its capacity
\citep[Asia service-provider disclosure]{jpmprivatebank2026}.

The US Treasury's Office of Foreign Assets Control (OFAC) includes foreign
branches of US entities among US persons. A separately incorporated foreign
subsidiary requires a program-specific inquiry: some programs extend obligations
to subsidiaries owned or controlled by US persons. The system must preserve
that relationship and apply the particular program's definition
\citep[FAQ 11]{ofacuspersons2024}. For the Chinese securities regime, Executive
Order 13959, as amended by Executive Order 14032, includes foreign branches
within its US-entity definition, alongside US citizens, lawful permanent
residents and persons in the United States. The investor's and the bank's
statuses must therefore be recorded separately
\citep[replacement section 3(d)]{eo140322021}.

The account's booking entity is only the beginning. Advice, execution, custody
and payment processing may involve different entities. Each participant needs
an identified role, location and applicable obligations. A decision about an
advisory service cannot automatically authorize principal dealing by the bank,
and permission for a purchase cannot automatically authorize a subsequent
transfer or conversion. These are distinctions the proposed data model must
retain even when a particular legal rule ultimately treats two cases alike.

\subsection{A US intermediary is different from a US investor}

The Chinese military-industrial securities restrictions supply a useful
counterexample to an overly broad institutional rule. Suppose, conditionally,
that the issuer is designated, the security is covered, the restriction has
taken effect, and no relevant authorization or divestment exception applies.
A purchase by a US person is then prohibited under the amended executive order.
Neither Chinese state ownership alone nor an instrument's AT1 label establishes
these premises; designation, security coverage and dates require their own
evidence \citep[replacement section 1]{eo140322021}.

Now change the bank's capacity while preserving those instrument facts. OFAC
permits specified investment advisory and management services for a non-US
person when the underlying transaction is otherwise permissible. Its guidance
explicitly addresses the ultimate benefit of US persons and evasion. This means
that a US bank's involvement does not by itself establish that every non-US
client transaction is prohibited \citep[FAQ 902]{ofaccmicservices2021}. OFAC also
states that US financial institutions may continue to intermediate qualifying
transactions between non-US persons. Conversely, a prohibited transaction under
this particular order must be rejected; the order does not itself require an
asset freeze \citep[FAQ 1048]{ofaccmicintermediation2022}. The system therefore
needs the actual service and beneficial parties, rather than an undifferentiated
``US bank involved'' flag.

Time and action introduce another distinction. OFAC permits continued holding
after the relevant divestment period, although a purchase or sale after that
period requires authorization. An instruction to sell an existing holding can
therefore require a different answer from an instruction to retain it
\citep[FAQ 1046]{ofaccmicholding2022}. The exact prohibition and divestment dates
depend on the designation history; an old prospectus cannot supply later list
changes \citep[FAQ 899]{ofaccmicdates2021}. Automatic conversions, redemptions
and other corporate actions require their own analysis. The present conditional
implementation supports buying, selling and holding; it leaves an unsupported
action unresolved.

Ownership also has program-specific meaning. Under OFAC's blocking ownership
rule, suppose blocked person X owns half of company A and A owns half of company
B. A becomes blocked through X's ownership, and B becomes blocked through A's
ownership. Multiplying the two fractions and treating B as merely one-quarter
owned by X misses the applicable propagation. Multiple blocked owners' direct
stakes must also be aggregated \citep[FAQ 401, examples 1--3]{ofacownership2014}.
Under the Chinese military-industrial securities regime, however, a subsidiary
must itself be designated: the blocking ownership rule does not apply merely
because its parent appears solely on the Chinese military-industrial securities list
\citep[FAQ 857]{ofaccmicsubs2021}. A reusable ownership graph consequently needs
a rule that states which relationships have legal consequences in each program.

\subsection{Compliance includes institutional and account obligations}

Sanctions are one source of additional restrictions. The retained regulation of
the Office of the Comptroller of the Currency (OCC) requires national banks to maintain a written Bank Secrecy Act
compliance program, with internal controls, independent testing, designated
responsibility and personnel training. It also addresses customer identification
within that program \citep[12 CFR 21.21(c)--(d), 2025 edition]{occbsa2025}.
These obligations concern the institution's conduct and controls. Reading a
security's prospectus cannot establish that the bank has performed them.

Applicability still precedes assessment. In the retained 2025 rules,
31 CFR 1010.620(a) addresses defined private-banking accounts established,
maintained, administered or managed in the United States. The account definition
adds conditions concerning the required minimum aggregate assets, non-US owners
and a liaison with the institution. Thus neither the commercial description
``private banking'' nor a Hong Kong office address alone decides whether that
particular US provision applies. Its substantive requirements include ownership,
source-of-funds and account-activity inquiries, with enhanced scrutiny for the
specified foreign political figures
\citep[31 CFR 1010.605(m), 2025 edition]{fincenprivatebankdef2025}
\citep[31 CFR 1010.620(a)--(d), 2025 edition]{fincenprivatebankdd2025}.
Those editions illustrate the scope distinction; their preservation does not
establish that every provision remains current for a later transaction.

The proposed system separates these institutional and account requirements from
product distribution and conduct rules. It also reserves a distinct place for
the bank's internal policies and the client's mandate. An internal restricted
list or product policy can prevent a service even where the examined public law
does not prohibit it. The repository has not received JPMorgan's actual internal
policies, so their requirements and permitted exceptions remain missing inputs.
They must not be reconstructed from a public marketing disclosure.

Finally, execution and custody require the actual settlement route and relevant
reporting duties. A requirement to investigate an account, reject a trade,
freeze specified property or make a report has a different legal effect. The
decision must preserve those distinctions. A generic negative eligibility flag
cannot safely stand for all four actions, and satisfying one duty does not
establish that all the others have been performed.

\subsection{Compose the obligations without erasing qualifications}

The common front end now needs to supply two kinds of information together:
the instrument's sourced terms and the transaction's parties, institutions,
services and dates. Rule packs selected for that context supply the applicable
requirements. In the Python implementation these are organized into seven
assessment categories: scope and coverage; sanctions; bank regulation and
licensing; client and account controls; product distribution and conduct; bank
policy and mandate; and settlement, custody and reporting. The categories are
a software design choice. They do not assert that every relevant rule has
already been discovered or implemented.

Each required assessment carries a result, a reason, source references and a
recorded validity interval. A finding that a rule does not apply needs evidence
just as a finding that its requirements are satisfied does. Unknown, conflicting,
expired and missing findings prevent an affirmative result. The binding also
identifies the instrument, client, booking entity, establishment, service,
action, transaction time, knowledge time, and content hashes for the facts,
settlement route and policy snapshot. Changing any of them invalidates the
earlier bound finding. Checking a source's hash establishes unchanged bytes;
the scope assessment must separately address legal currentness and meaning.

The conjunction has a simple conditional justification. Let $R$ be the finite
set of declared requirements. Let $S$ contain requirements supported as satisfied
or inapplicable by valid, consistent findings for this context. Let $U$ contain
requirements with unresolved findings, and let $C$ mean that every required
assessment category is represented. The implementation's permission within
that declared scope is
\[
   P = C \;\land\; (R\subseteq S) \;\land\; (U=\varnothing).
\]
A valid prohibition excludes its requirement from $S$, and a conflicting
favorable finding also places that requirement in $U$. Either prevents $P$.
An absent bank-policy assessment has the same practical consequence for
permission, but its explanation is missing evidence rather than a proved legal
prohibition. A favorable product assessment cannot alter either result.

Representing every category still leaves a possible omission. A program could
receive one favorable assessment in each category while never being asked
about another applicable restriction. The integrated investigation therefore
starts from a versioned inventory of fourteen declared obligations. A request
cannot shorten that inventory or supply its own completed findings. Requirements
without an executable interpretation remain visible, including other sanctions
programs, licensing, issuer law and the unavailable internal bank policy. This
prevents omissions within the declared inventory; discovering every obligation
outside it remains a separate legal question.

The factual references now retrieve the underlying documents. For a structured
account record, the program reads the named field, checks its type and preserves
disagreement between records. Source dependencies, effective dates, knowledge
dates and amendments are checked separately. Consider the private-account
definition above: the required minimum deposit and the places where the account
is administered or managed enter the calculation as distinct premises. A large
current balance cannot replace the required contractual minimum. A quoted
provision or an unchanged document still does not establish that a proposed
interpretation follows from it.

The resulting conditional specifications use the shared RuleIR and Catala
translation machinery. They include entity scope, selected sanctions,
institutional controls, private-account diligence and reporting requirements,
together with the existing product calculations. Formal-expression preservation
is checked separately from compiled execution. An independently written
mathematical specification supplies another comparison for the Boolean rules
and exact monetary thresholds. Agreement between programs therefore supports
their stated formal consequences while leaving the interpretation of the legal
sources explicit.

Changes after an assessment also matter. The event record distinguishes a new
source, an account change, a different settlement route and a corporate action.
Each requires reassessment before the earlier decision can be reused. Reporting
and other duties retain their responsible party, trigger, deadline and supporting
evidence. A recorded completion document establishes that evidence was supplied;
it does not independently establish that the real action occurred. When the
integrated example uses the preserved public documents without a real account
or internal bank policy, its transaction decision consequently remains qualified.

The checks enumerate every combination of satisfied, prohibited and unknown
across the seven categories, giving $3^7=2{,}187$ cases. Further cases challenge
conflicts, non-applicability, missing policy, source changes and changed entity,
action or time. The ownership calculation has an independent finite check: for
each of 512 three-entity graph and initial-blocking combinations, the test
intersects every ownership-closed superset of the initial blocked entities and
compares that set with the implementation's result. The algorithm is monotone
and terminates on a finite graph; induction shows that every closed superset
contains each entity it adds. Its terminal set is therefore the least closed
superset under the supplied ownership premises.

No human quality labels are used in these checks. They establish properties of
the declared formal rules and the handling of missing evidence. Correctness of
a new legal interpretation, completeness of an ownership register and coverage
of future amendments remain separate obligations. The system can generalize by
applying checked semantics to new supported facts and qualifying unsupported
situations; it cannot turn today's finite legal collection into a proof about
every law that might govern a future transaction.

\subsection{One investor, two CoCo offering decisions}
\label{sec:coco-purchaser-cases}

A useful positive and negative test can hold the investor fixed while changing
the security. Consider a hypothetical individual resident and located in the
United States, buying USD~200,000 principal for their own account and benefit.
Compare the original offering terms of two actual contingent convertible
securities: Barclays PLC's February 2025 issue and Standard Chartered PLC's
January 2025 issue. The question is whether that investor satisfies the
specified US purchaser conditions. Approval below means passing this particular
condition; it does not mean that a regulator approved the investment or that
the bank may execute the transaction. The different issue dates are retained:
this is a comparison of preserved offering terms, not an assertion that both
primary offers are open today.

The relevant investor category needs explanation before either document can be
applied. A \emph{qualified institutional buyer}, or QIB, is a category defined
by Rule~144A for specified institutions and accounts. An individual's wealth
or Hong Kong professional-investor status does not turn that individual into
a QIB. For this case, non-QIB status is an explicit hypothetical premise.
Regulation~S uses a separate definition of US person, which includes a natural
person resident in the United States. The branch and account exceptions in that
definition also show why the sanctions definition of US person discussed above
must not simply be reused. The comparison retains both the purchaser's status
and the beneficiary of the purchase
\citep[Rule~144A(a)(1); Regulation~S, section~230.902(k)]{rule144acoco2025,regsdefinitions2025}.

The first security is Barclays PLC's USD~1.5~billion 7.625\% Fixed Rate Resetting
Perpetual Subordinated Contingent Convertible Securities, with international
securities identification number (ISIN)
\texttt{US06738EDC66}. The supplement is dated 18~February 2025 and identifies
25~February as the expected delivery date. It identifies its Form~F-3
registration statement; the underwriting section describes an offer directly
to the public. The minimum denomination is USD~200,000, with USD~1,000 increments.
These affirmative offering provisions support the declared public US
distribution route. In particular, ``registered form'', which describes how
ownership is recorded, is not itself the evidence of Securities Act registration
\citep[cover; pp.~S-6, S-22 and S-99]{barclayscoco2025}.

The second security is Standard Chartered PLC's USD~1~billion 7.625\% Fixed Rate
Resetting Perpetual Subordinated Contingent Convertible Securities. Its offering
circular is dated 8~January 2025 and specifies issue on 16~January. The restricted
and offshore identifiers are \texttt{US853254DF47} and \texttt{USG84228GP72}.
It has the same minimum denomination and increment. Its US distribution,
however, is limited to QIBs; the alternative described in the circular is an
offshore distribution to non-US persons. The purchaser representations also
distinguish acquisition for one's own account from acquisition for a QIB's
account. Our individual satisfies neither the institutional route nor the
offshore route \citep[pp.~i, 14--15, 140 and 148--149]{standardcharteredcoco2025}.

The order size therefore does not explain the difference. Under the declared
interpretations, the comparison is:
\begin{center}
\begin{tabularx}{\linewidth}{@{}lXl@{}}
\toprule
Original issue & Distribution relevant to this purchaser & Purchaser condition \\
\midrule
Barclays, February 2025 & Public US offering & Pass \\
Standard Chartered, January 2025 & QIB-only US distribution & Fail \\
\bottomrule
\end{tabularx}
\end{center}
Neither row establishes SFC product approval or an entitlement to an allocation
from the underwriters.

The conditional rule can now be stated precisely. Let $R$ mean that the
declared route is the public US distribution, and $T$ that it is the specified
Rule~144A/Regulation~S distribution. Let $Q$ mean that the purchaser is a QIB
and buys for its own account or a QIB's account. Finally, let $O$ mean that the
transaction is offshore, the purchaser and beneficial account satisfy the
specified non-US conditions, and the purchaser is not an issuer affiliate.
The selected purchaser condition is
\[
   G = R \;\lor\; \bigl[T\land(Q\lor O)\bigr].
\]
For Barclays, the declared public distribution makes $R$ true. For Standard
Chartered, $R$ is false and $T$ is true, but our investor makes both $Q$ and $O$
false. The equation therefore yields opposite results from the same investor
premises. It represents these purchaser conditions, not every condition of
either securities-law exemption, every jurisdiction's selling restrictions or
the effectiveness of a registration statement.

Changing the purchaser tests whether this result has been tied incorrectly to
an issuer's name. A QIB buying for its own account changes the Standard
Chartered result to a pass under the same declared rule. A bank that is a QIB
but buys for a non-QIB customer's account does not satisfy $Q$ merely through
the bank's status. Moving the original individual offshore also leaves a
problem: US-person status and US beneficial ownership do not disappear with
the location of execution. These changes distinguish purchaser, account and
geography without supplying an answer label for either issuer.

The loss mechanisms remain relevant to the other assessments. Both documents
describe contingent conversion and regulatory powers that can reduce or
extinguish an investment. Passing a purchaser condition is consequently no
evidence of principal protection. Hong Kong distribution, suitability,
sanctions, bank authority, internal policy and settlement still require their
own conclusions. The US purchaser comparison cannot replace those conclusions
\citep[Barclays cover and pp.~S-16, S-102; Standard Chartered pp.~i, 140--142]{barclayscoco2025,standardcharteredcoco2025}.

The executable comparison reproduces the two results through both RuleIR and
Catala. Its 157 input cases produce 314 executions: each backend computes the
independently derived result for 138 complete cases and abstains on 19 cases
with missing, conflicting or expired evidence. A separate mathematical check
compares the implemented expression with an independently written relation
over all $2^9=512$ assignments of its nine Boolean premises. A proof checks
preservation of the declared formal expression in the shared model, and six
deliberately wrong variants are detected, including omission of the beneficial
account condition. The focused regression suite passes 165 tests. These results
establish conditional formal and software properties. The legal interpretation
and completeness of the applicable requirements remain unproved.

In the integrated bank investigation, the Barclays purchaser result remains
qualified by missing transaction evidence. The Standard Chartered result
additionally prevents proceeding under the declared offering route. Both retain
all seven assessment categories and the fourteen declared bank obligations;
neither permits execution of an actual trade. The preserved prospectuses,
explicit interpretation proposals, synthetic investor premises and reproducible
checks make this a useful test without human \mbox{quality labels}. New issuers can
reuse the checked relation when its premises and scope apply; new offering
terms or unsupported legal situations require a new interpretation or a
qualified answer. The finite pair cannot establish future legal correctness.

\subsection{Connecting the loss provisions to the selling rules}
\label{sec:joined-eligibility}

The next calculation follows the same investment from its loss mechanism to
the Hong Kong selling requirements. The provisions surveyed in
Section~\ref{sec:coco-legal-risk} supply the questions: what legal form does the
instrument take, which contingent loss mechanism applies, which institution
provides the service, and what is known about the purchaser? A contractual
write-down, a conversion and a regulator's power are separate facts. The
retained UBS notes provide for conversion. Cancelling their debt principal
would not, by itself, establish that the shares received are worth zero.

The implementation connects seven conditional calculations. It considers direct
debt with the specified contingent loss features and separately records whether
an investment product falls within the related-product provisions. An
equity-form preferred share does not become direct debt merely because it can
lose value. Related-product scope remains document-dependent; the program does
not invent a percentage threshold for a fund's principal strategy. It then
examines the intermediary arrangement, the PI restriction, the specified PI
exemptions, the resulting duties, the issuer's legal powers and any exercise of
those powers. Each conclusion retains its required premises.

The execution-broker exception shows why this decomposition matters. The
earlier calculation accepted an established exception as one premise. The new
calculation expands it into the specific adviser or asset-manager arrangement,
an order placed through that arrangement, the adviser's regulatory status,
execution and custody services only, absence of direct customer contact, the
written allocation of responsibility, written customer notification and the
continuing arrangement. A direct customer request and circumvention each
prevent reliance on the exception. These are proposed formal consequences of
the retained FAQ, not evidence that a customer's arrangements meet its
conditions. A bare assertion that the exception applies cannot bypass an
unknown condition or contradict a known one.

Duties are computed separately from the selling restriction. The specified
institutional and corporate PI exemptions affect risk rating, enhanced
suitability and disclosure; the PI selling restriction remains separate.
Loss-absorption funds receive the specified treatment for enhanced suitability
and disclosure, with other SFC requirements retained. Discretionary portfolio
management can change the level at which an assessment is performed. It does
not supply evidence that the assessment has been performed. Additional
disclosure flexibility referred to by the FAQ still needs its own applicable
authority and conditions.

A regulator's power requires a further chain. The jurisdiction, entity,
instrument, legal basis, triggering circumstances and procedure must support
the particular power. An applicable order must then have been issued, its
suspension status matters, and implementation is another event. The calculation
therefore distinguishes authority from an operative order and an implemented
loss. It cannot infer any of them from the word AT1. For the UK examples, the
joined report retains a missing instrument-specific UK resolution basis. For
UBS, it retains the dated Swiss source and its currentness qualification.

Contractual amounts use exact rational arithmetic under explicitly supplied
event conditions. A hypothetical conversion of principal 3,777 at a
same-currency conversion price of 37.77 gives 100 shares and zero remaining
debt principal. The market value of those shares is a different quantity.
If notice, timing or other instrument-specific event conditions are unknown,
the calculation remains unknown. Section~\ref{sec:ubs-terms} develops the
instrument-specific notice and share calculations beyond this generic example,
including the price branches that still prevent a unique result.

The joined result recomputes the bank assessment and retains all fourteen
obligations in its seven categories. It compares the shared PI premise across
the product and bank inputs, preserving contradictory observations. A changed
source, fact document, mandate or execution route requires reassessment.
These calculations supplement the broader product-law obligation; they do not
mark it complete. With the retained public documents and no actual account,
mandate or transaction evidence, the UBS, Barclays and Standard Chartered
reports remain qualified and cannot authorize execution.
